%% =====================================================================
\section{Introduction}\label{sec:intro}
%% =====================================================================

\subsection{The three known measures}\label{ss:three}

Let $p$ be a prime and let $\zeta_p$ denote the Kubota--Leopoldt $p$-adic zeta function,
normalised as in \cite{LSZ2025} so that $\zeta_p(s)=L_p(s,\omega^{1-s})$ for integers
$s\ge2$. Only three $p$-adic irrationality \emph{measures} are on record. Writing
$\mu(\xi)$ for the irrationality exponent of $\xi\in\Qp\setminus\QQ$ in the $p$-adic sense
of Bel \cite{Bel2019}, they are
\begin{equation}\label{eq:three}
\begin{aligned}
 \mu(\zeta_2(3)) &\le \frac{12\log2}{6\log2-3} = 7.177398\dots,\\
 \mu(\zeta_3(3)) &\le \frac{6\log3}{3\log3-3} = 22.281447\dots,\\
 \mu(\zeta_2(5)) &\le \frac{16\log2}{8\log2-5} = 20.342651\dots .
\end{aligned}
\end{equation}
The first two are extracted in \cite[\S1]{LSZ2025} from Calegari's work \cite{Calegari2005}
by way of Bel's lemma, or equivalently from Beukers \cite[Thm.~11.2]{Beukers2008}; the third
is the main theorem of \cite{LSZ2025}. The corresponding irrationality statements go back to
Calegari \cite{Calegari2005} and Beukers \cite{Beukers2008} for weight $3$, and to
\cite{LSZ2025} for $\zeta_2(5)$; the many-values theorems of Lai--Sprang
\cite{LaiSprang2023} and Lai \cite{Lai2024} produce irrationality of \emph{one of} a list,
never a measure for an individual value.

Every one of the constructions behind \eqref{eq:three} is \emph{totally symmetric}: the
underlying rational function has all its numerator bricks of the same length and all its
pole bricks of the same length, sitting on the same interval. In the archimedean world that
is precisely the configuration on which the Rhin--Viola arithmetic group method
\cite{RhinViola}, and its hypergeometric reformulation by Zudilin \cite{Zudilin2004}, buy
\emph{nothing}: the orbit of a totally symmetric point is a single point
(Proposition~\ref{prop:symzero}). The archimedean record $\mu(\zeta(3))\le5.51389062$ of
Rhin and Viola lives instead at a strongly asymmetric point, where the group removes about
$40\%$ of the denominator (\S\ref{ss:howmuch}).

\subsection{The questions we answer}\label{ss:questions}

Lai, Sprang and Zudilin end \cite{LSZ2025} with two explicit invitations. The first
\cite[Final remarks]{LSZ2025}:

\begin{quote}
``Speaking about the irrationality measure for $\zeta_2(5)$ and also for the related
$p$-adic zeta values $\zeta_2(3),\zeta_3(3)$, our construction in this note and the earlier
ones in [Beukers, Calegari, Lai] exploit exclusively the so-called totally symmetric
hypergeometric approximations (or their equivalents). It may certainly be of interest to
take a step further and investigate more general hypergeometric series leading to rational
approximations that depend not only on $n\in\ZZ_{\ge0}$ but on multiple integer parameters.
Such generalisations are amenable to additional arithmetic manipulations, for example, to
the use of the arithmetic group method as developed in the works of Rhin and Viola.''
\end{quote}

\noindent The second, after observing that Lai's and Beukers' two rational functions produce
the \emph{same} linear forms in $1$ and $\zeta_2(3)$ \cite[Final remarks]{LSZ2025}:

\begin{quote}
``It is natural to expect a hypergeometric-type identity for the Volkenborn integrals behind
the coincidence $S_n^{(\mathrm L)}=S_n^{(\mathrm B)}$.''
\end{quote}

\noindent Independently, Brown \cite[Rem.~47]{Brown2026} records a referee's question in the
%% REFEREE [R13]: the quotation was truncated mid-sentence, silently promoting Brown's
%% comma to a full stop and dropping the clause about prime cancellation already visible
%% in his experiments -- which is material to \S\ref{ss:evade}(A).  Restored in full.
same spirit for his determinant criterion: ``A referee asked the excellent question of
whether the matrices $D_n$ can be designed to make use of the cancellation of primes in the
method of Rhin--Viola. I do not know the answer to this question, but the numerical
experiments to be discussed later clearly demonstrate a large degree of `prime
cancellation' of possibly a different nature.''

This paper answers the second question completely and the first question
quantitatively --- the latter in the negative, but in a way that we believe is more useful
than a positive answer of the expected size would have been, because it comes with a map.

\subsection{Results}\label{ss:results}

\smallskip\noindent\textbf{(A) The Transfer Principle} (Theorem~\ref{thm:transfer},
\PROVED). Let $R\in\QQ(t)$ have poles only at non-positive integers and $\deg R\le-2$, with
partial fractions $R(t)=\sum_{i,k}r_{i,k}/(t+k)^i$, and set
\[
 \rho_i:=\sum_k r_{i,k},\qquad
 \rho_{0,\theta}:=-\sum_{i,k}r_{i,k}\sum_{\nu=0}^{k-1}(\nu+\theta)^{-i}.
\]
Then, with $\widetilde R$ a primitive of $R$,
\[
 \sum_{m\ge0}R(m+\theta)=\rho_{0,\theta}+\sum_{i\ge2}\rho_i\,\zeta(i,\theta),
 \qquad
 -\int_{\Zp}\widetilde R(t+\theta)\,\dd t
   =\rho_{0,\theta}+\sum_{i\ge2}\rho_i\,\omega(\theta)^{1-i}\zeta_p(i,\theta),
\]
\emph{with the same rational numbers} $\rho_{0,\theta},\rho_2,\rho_3,\dots$. The two
constructions differ only in which \emph{value} is attached to each $\rho_i$; the
coefficient map $R\mapsto(\rho_{0,\theta},\rho_2,\rho_3,\dots)$ is literally the same
$\QQ$-linear functional of the partial-fraction data. The immediate corollary
(Corollary~\ref{cor:transfer}) is that every Rhin--Viola/Zudilin transfer identity, being an
identity between the rational coefficients at two parameter points, is simultaneously an
archimedean and a $p$-adic identity; likewise every denominator and integrality lemma. In
particular the second LSZ question is answered by observing that no hypergeometric-type
identity \emph{for the Volkenborn integrals} is needed: the identity lives on the rational
coefficients, and the Volkenborn integral never sees it.

\smallskip\noindent\textbf{(B) The group, made explicit and made $p$-adic}
(\S\ref{sec:group}). For the weight-$3$ family we realise Zudilin's group
$\Gp=\langle\mathfrak a_1,\mathfrak a_2,\mathfrak a_3,\mathfrak b,\mathfrak h\rangle$ of
order $1920=|W(D_5)|$ \VER{BFS closure, exact}, exhibit its invariant in the
$h$-parametrisation, and verify the transfer identity coefficientwise --- on both the
rational part and the $\zeta(3)$-part --- at every point of the orbits of four directions,
$1200/1200$ at the generic direction with $0$ failures \VER{4 directions, orbits of size
$1,40,80,1920$}. Crucially for the $p$-adic application, the same test passes at
\emph{half-integer} parameters, $780/780$ comparable points with $0$ failures: the mechanism
is lattice-robust, not an accident of integrality. The savings function $\dgr$ is
implemented and validated end to end by recovering Rhin and Viola's published
$\mu(\zeta(3))\le5.51389062$ as $5.51389063$ \VER{$8$ digits}. A cross-tabulated sweep of a
$37$-paper corpus for $p$-adic markers against group markers returns no paper using a
Rhin--Viola group $p$-adically \VER{corpus of $37$ papers, \S\ref{ss:corpus}}.

\smallskip\noindent\textbf{(C) The ledger and its validation} (\S\ref{sec:ledger},
\S\ref{sec:validation}). We write down a parametrised Volkenborn construction containing, as
exact special cases with their printed constants, the constructions of Beukers, Calegari
(through Beukers' reformulation), Lai, Lai--Sprang, Lai--Lupu--Sprang and LSZ, and derive an
exact cost ledger
\[
\begin{gathered}
 G=v_p(C)\log p,\qquad
 \alpha=\Bigl[v_p(C)+Ar+\min_{\theta_0\in\Theta}\sum_c\ctb_c(\theta_0)\Bigr]\log p,\\
 \beta=G+E,\qquad \marg=\alpha-\beta ,
\end{gathered}
\]
in which every component is separately derived and separately measured. Fed \emph{only} the
parameter tuple, the ledger returns all three measures \eqref{eq:three} to every printed
digit, plus four further independent anchors. This validation section is the load-bearing
part of the paper's credibility: nothing downstream is calibrated.

\smallskip\noindent\textbf{(D) The classification} (Theorem~\ref{thm:classification},
\PROVED). A single Volkenborn form in $1$ and the \emph{full} value $\zeta_p(w)$, rather
than in Hurwitz values $\zeta_p(w,\theta)$, exists exactly when $\varphi(p^r)=2$, i.e.\
$p^r\in\{3,4\}$, plus the exceptional point $(p,r)=(2,1)$ where the single shift $\tfrac12$
already \emph{is} the full $\zeta_2$. Otherwise the full twisted coset sum is forced, and
the $p$-adic smallness becomes a minimum over $\varphi(p^r)$ cosets. At $p\ge5$, whenever
$\varphi(p^r)$ exceeds the number of distinct numerator-brick shifts, that
minimum is attained at a coset with no aligned brick, and then
$Ar+\sum_c\ctb_c=0$ \emph{exactly}: the $p$-adic smallness cancels the archimedean growth
to the last digit and the entire $d_n^E$ denominator cost goes unpaid, so $\marg=-E$
(Proposition~\ref{prop:pge5}). Without that hypothesis the deficit is bounded rather than
equal to $-E$ (Corollary~\ref{cor:noAhelp}, Remark~\ref{rem:pge5sharp}).
%% REFEREE [R1]: this passage asserted margin = -E for every rank-1 configuration at
%% p >= 5.  False: at (p,r,A,m,eps) = (5,1,4,1,1) with the four bricks on the four
%% cosets j/5 the worst coset still carries one aligned brick, contributions
%% (-1,-1,-1,+1/4), so A r + sum ctb = 1.25 and margin = 1.25 log 5 - E = -2.9882.
%% That is the (p,w) = (5,5) cell of Finding~\ref{find:map1}.  Hypothesis restored.

\smallskip\noindent\textbf{(E) Optimality of the three known measures}
(Theorem~\ref{thm:optimal}, \COMP{stated cone}). Over the cone described precisely in
\S\ref{ss:optdomain} --- all $O(1)$ insets of the bricks, all scaled (Rhin--Viola) insets,
all rank reductions by polynomial-coefficient difference operators, and the whole
$1920$-element group orbit --- the bounds \eqref{eq:three} are attained and cannot be
improved:
\[
 \mu(\zeta_2(3))\le7.17739889912418,\quad
 \mu(\zeta_3(3))\le22.28144795149432,\quad
 \mu(\zeta_2(5))\le20.34265173891448,
\]
with $E$ and the archimedean growth \emph{measured} rather than modelled, so that any
$\Phi_n$-type or Rhin--Viola saving actually present is already counted. As stressed in
\S\ref{ss:scope}, this is optimality over a family, not an absolute statement.

\smallskip\noindent\textbf{(F) The rank wall} (\S\ref{ss:rank}). For the very-well-poised
$p=2$ family a parity law forces $\rho_i=0$ whenever $i+M$ is even
(Proposition~\ref{prop:parity}), whence the whole ledger collapses to
$\marg(M,m)=2M\log2-(M+m)$, which tends to $+\infty$: \emph{size never binds at $p=2$}. The
sole obstruction is that the form carries $\lfloor(M-1)/2\rfloor$ zeta values. Removing a
coefficient requires a difference operator with polynomial coefficients; five searches over
inset cones (dimensions $144,210,144,144,112$) return $0$ out of $300$ kernel vectors that
keep the wanted coefficient, against positive controls at $300/300$
\VER{$\FF_q$, $q=2^{61}-1$, every hit re-verified over $\QQ$}.

\smallskip\noindent\textbf{(G) The shut hatch} (\S\ref{ss:hatch}). The only place where the
$p\ge5$ verdict could be overturned inside this parameter space is cancellation between
cosets, since $\alpha=\min_j$ is an inequality. Exact computation of every $S_n(j/p^r)$
shows the plain and all $\omega$-twisted sums beat the per-coset minimum by at most $3$
\emph{absolutely}, not per $n$; and the Smith-normal-form hyperplane depth of the coset
vectors is bounded by $28$ and constant in $n$, while the requirement grows like $E\cdot n$
\VER{$p\in\{5,7\}$, $r\in\{1,2\}$, $A\in\{2,3,4\}$, $m\in\{0,1\}$, $n\le30$}.

\smallskip\noindent\textbf{(H) The map and the dichotomy} (\S\ref{sec:map}). Putting the
above together gives a complete margin table over $p\le13$, $w\le9$ and rank $\le4$, in two
versions: with the ledger's $E$, and with the corrected, measured $\Etrue$. The headline is
a dichotomy that the campaign literature has been treating as one obstruction:
\begin{center}
\emph{at $p\in\{2,3\}$ the binding wall is rank; at $p\ge5$ it is size.}
\end{center}
After the $E$-correction the nearest misses in the whole family are $\zeta_5(3)$ and
$\zeta_7(3)$ at exactly $-3.000$, and those are structurally the \emph{hardest} cells, since
there the deficit is $-E$ exactly.

\smallskip\noindent\textbf{(I) The structural finding} (\S\ref{ss:tension}). Bel's $p$-adic
criterion requires $\alpha>\text{growth}+\text{denominators}$, whereas the archimedean
Rhin--Viola criterion requires only $C_0>C_2$: the archimedean coefficient growth $C_1$
enters the archimedean \emph{measure} but never the archimedean \emph{irrationality
condition}, while $p$-adically it must be paid. And $C_1=0$ holds exactly on the degenerate
locus where the numerator bricks overlie the poles --- which is exactly where the group
saving $\dgr$ vanishes. \emph{$\dgr>0$ and $C_1=0$ have disjoint support.} Measured,
$C_1/\bud$ is $1.175$, $0.969$, $1.073$ at three test directions, each alone exceeding the
entire available $\dgr/\bud\le0.44$.

\subsection{Scope of the optimality claims}\label{ss:scope}

We state this once, prominently, and repeat it at each use. \emph{No claim in this paper is
that the bounds \eqref{eq:three} are unimprovable.} What we prove and compute is optimality
\emph{over the family and cone that we search}, which is stated explicitly, with the search
domain written out, at every occurrence. The family is the parametrised Volkenborn
construction of Definition~\ref{def:family} together with the inset cones of
\S\ref{ss:optdomain} and the difference-operator module of \S\ref{ss:rank}. A construction
outside that family is not covered; \S\ref{ss:evade} describes what such a construction
would have to look like, precisely so that the boundary of our claim is visible.

Similarly, the negative results of \S\ref{ss:rank} and \S\ref{ss:hatch} are fair sweeps with
clean negatives and stated ranges, backed by positive controls; they are not impossibility
proofs. \S\ref{ss:openlemmas} isolates the two lemmas that would convert them into theorems.

\subsection{Evidence classes}\label{ss:evidence}

Every substantive assertion below carries one of the following tags. This convention is
inherited from the working memos underlying this paper and is used without exception.

\begin{itemize}[leftmargin=2.2em,itemsep=2pt]
\item \PROVED{} --- a proof is given here, or read out of a proof in a cited source, with the
      source lemma named.
\item \DERIVED{} --- derived here from cited lemmas by a chain that is written out, but not
      separately measured.
\item \VER{range} --- an exact finite computation in exact rational or integer arithmetic; the
      range of parameters and the sample size are always stated inside the tag.
\item \COMP{domain} --- a computed optimum, exhaustive over the stated domain; the domain is
      always stated inside the tag.
\item \EXCL{tolerance, bounds} --- a search that returned nothing, with the tolerance and the
      coefficient/degree bounds of the search stated.
\item \CALIB{} --- a modelling assumption. Exactly one such statement appears in this paper
      (\S\ref{ss:calibrated}), it is flagged in place, and it was subsequently
      \emph{refuted} by the first-principles computation of \S\ref{ss:tension}; we retain it
      because the refutation is the structural finding.
\item \OPENC{} --- a statement we believe and do not prove.
\end{itemize}

All arithmetic reported as exact is carried out in \texttt{fractions.Fraction} or in exact
integer arithmetic; $p$-adic valuations come from exact truncated Volkenborn series.
Wide combinatorial searches are run over $\FF_q$ with $q=2^{61}-1$ and every hit is
re-verified over $\QQ$. Floating point appears only in the final evaluation of
transcendental constants such as $\log2$, and the headline constants are re-verified at
$50$ decimal digits.

\begin{remark}[Methodology; flagged]\label{rem:method}
The computations reported here were carried out with substantial machine assistance,
including automated search over parameter cones and over candidate difference operators;
every positive identification was re-derived independently of the search that proposed it,
every negative is reported with its search domain, tolerance and positive control, and the
programs are listed in Appendix~\ref{app:repro}.
\end{remark}

\subsection{Notation and one renaming}\label{ss:notation}

$p$ is a prime, $q_p=p$ for odd $p$ and $q_2=4$, $\varphi$ is Euler's totient, $v_p$ is the
$p$-adic valuation with $v_p(p)=1$, and $d_n=\lcm(1,\dots,n)$. We write $\omega$ for the
Teichm\"uller character $\ZZ_p^\times\to\mu_{\varphi(q_p)}(\ZZ_p)$, \emph{extended to}
$\Qp^\times$ as in \cite[\S2]{LaiSprang2023} by
\begin{equation}\label{eq:omegaext}
 \omega(x):=p^{v_p(x)}\,\omega\bigl(x/p^{v_p(x)}\bigr),\qquad
 \langle x\rangle:=x/\omega(x).
\end{equation}
%% REFEREE [R3]: only the Teichmuller character on Z_p^times was defined, but
%% omega(theta)^{1-i} in \eqref{eq:padic} and omega(a)^{-u} in \eqref{eq:Sn} are applied
%% at |theta|_p > 1.  The extension is load-bearing: it carries the factor p^{r(i-1)}
%% that reconciles Theorem~\ref{thm:transfer}(ii) with Proposition~\ref{prop:linform}.
The extension is what carries the powers of $p$ in \eqref{eq:padic} and \eqref{eq:Sn}. We write $(x)_L=x(x+1)\cdots(x+L-1)$ for the
Pochhammer symbol and $\zeta(s,\theta)$, $\zeta_p(s,\theta)$ for the Hurwitz and $p$-adic
Hurwitz zeta functions. Following \cite[Lemma 12]{LaiSprang2023} we use
\begin{equation}\label{eq:cosetsum}
 D^i\,\zeta_p(i)=\sum_{\substack{1\le j\le D\\ \gcd(j,p)=1}}
   \omega\Bigl(\frac jD\Bigr)^{1-i}\zeta_p\Bigl(i,\frac jD\Bigr),
 \qquad q_p\mid D,\ i\ge2 .
\end{equation}

\begin{convention}[the well-poised flag]\label{conv:rename}
In the source memos the very-well-poised indicator of the rational function is written
$\delta\in\{0,1\}$, and the Rhin--Viola denominator saving is \emph{also} written $\delta$.
Since both appear in the same formulas here, we write $\epsilon\in\{0,1\}$ for the
well-poised flag (the exponent of the factor $2t+h_0$) and reserve
$\dgr=\lim_n\log\Phi_n/n$ for the group saving. Consequently the denominator cost reads
$E=A+m+1-\epsilon$ and never $A+m+1-\delta$.
\end{convention}

\subsection{Plan}\label{ss:plan}

\S\ref{sec:transfer} proves the Transfer Principle and its corollaries.
\S\ref{sec:group} makes the group explicit, records the verification of the transfer
identity at integer and at half-integer parameters, states the four honest flags
(F1)--(F4) including the Gauss-duplication bridge, and reports the corpus sweep.
\S\ref{sec:ledger} sets up the parametrised construction and the exact ledger.
\S\ref{sec:validation} is the validation section.
\S\ref{sec:classification} proves the $\varphi(p^r)=2$ classification and the $p\ge5$ exact
cancellation.
\S\ref{sec:scan} reports the scan: optimality, the rank wall, the shut hatch.
\S\ref{sec:map} gives the margin map and the dichotomy.
\S\ref{sec:discussion} is the discussion: the two open lemmas, the disjoint-support tension,
and what a construction evading the map would need. Appendix~\ref{app:repro} is the
reproduction record.
